\chapter{Convexité, tangentes et inégalités}
\section*{Exercice 15.1 -- Convexe mais décroissante}
[Base] Étudier variations et convexité de $e^{-x}$.

\section*{Exercice 15.2 -- Inflexion}
[Méthode] Déterminer les points d'inflexion de $f(x)=x^4-2x^2$.

\section*{Exercice 15.3 -- Tangente logarithmique}
[Preuve] Démontrer $\ln x\le x-1$ pour $x>0$, avec le cas d'égalité.

\section*{Exercice 15.4 -- Milieu}
[Preuve] Montrer $e^{(a+b)/2}\le(e^a+e^b)/2$.

\section*{Exercice 15.5 -- Un faux critère}
[Critique] Expliquer pourquoi $f''(0)=0$ ne suffit pas à produire une inflexion.

\section*{Exercice 15.6 -- Lorenz}
[Application] Pour les revenus fictifs $1,1,2,6$, construire les points d'une courbe de Lorenz et vérifier la convexité.
